\section{The integral lift: the precise bridge from CMs to \texorpdfstring{$J$}{J}}\label{sec:lift}
\subsection{A common source, not an embedding of characteristic two}
The successful construction changes the amplitude arithmetic. Retain the four atomic CM positions, but allow an integer number of contributions at each position:
\begin{equation}
 \widetilde A=n_0E_0+n_1E_1+n_2E_2+n_3E_3,
 \qquad n_j\in\Z.
\end{equation}
Positive counts can describe accumulated paths; negative counts are a useful reduced representation. Addition is now ordinary integer addition. With cyclic convolution these objects form $\Z[C_4]$. Reducing their coefficients modulo two recovers $\alg_2$, but this reduction discards multiplicity information.

Next impose an additional phase relation: opposite positions represent opposite amplitudes. In algebraic terms, quotient by $1+g^2$. The construction is the diagram
\begin{equation}\label{eq:commonlift}
 \begin{tikzcd}[column sep=large]
 & \Z[C_4] \arrow[dl,"\bmod 2"'] \arrow[dr,"g^2=-1"] & \\
 \F_2[C_4] && \Z[g]/(g^2+1)\cong\Z[i].
 \end{tikzcd}
\end{equation}
The opposite-phase relation is \emph{a new semantic choice}, not a theorem of Boolean truth tables. There is no unital ring homomorphism from the characteristic-two algebra on the left into the characteristic-zero algebra on the right: the equality $1+1=0$ would have to remain true. Consequently the diagram must not be replaced with a purported inclusion $\F_2[C_4]\subset\Z[i]$.

\subsection{The intertwining equation}
Define
\begin{equation}\label{eq:QJ}
 Q=\begin{pmatrix}1&0&-1&0\\0&1&0&-1\end{pmatrix},\qquad
 J=\begin{pmatrix}0&-1\\1&0\end{pmatrix}.
\end{equation}
Then $Q\widetilde A=(n_0-n_2,n_1-n_3)^T$. The four elementary images are
\begin{equation}
 QE_0=(1,0),\quad QE_1=(0,1),\quad QE_2=(-1,0),\quad QE_3=(0,-1).
\end{equation}
Here $QE_j$ abbreviates $Qv(E_j)$.

\begin{theorem}[CM rotation induces the quarter-turn]\label{thm:intertwiner}
The CM rotation matrix $P$ and the signed quarter-turn matrix $J$ satisfy
\begin{equation}\label{eq:intertwiner}
 QP=JQ.
\end{equation}
The kernel of $Q$ is the ideal generated by $1+g^2$, so $J$ is the action induced by rotation on the quotient $\Z[C_4]/(1+g^2)$.
\end{theorem}
\begin{proof}
For $n=(n_0,n_1,n_2,n_3)^T$, $Pn=(n_3,n_0,n_1,n_2)^T$, hence
\begin{equation}
 QPn=(n_3-n_1,n_0-n_2)^T=JQn.
\end{equation}
The kernel consists of $(a,b,a,b)$, corresponding to $(a+bg)(1+g^2)$. This is precisely the relation identifying $E_2$ with $-E_0$ and $E_3$ with $-E_1$.
\end{proof}
This is the missing bridge behind the initial analogy. It is stronger than observing that two transformations both have period four. It also makes the limitation unambiguous: $J$ is not one of the sixteen binary CMs. It is a two-dimensional signed representation \emph{induced from} their four-position rotation after lifting and quotienting.  The abstract quotient $\Z[C_4]/(1+g^2)\cong\Z[i]$ is standard algebra; the role of \cref{thm:intertwiner} is to pin that quotient to the CM coordinate convention and to expose exactly where the scalar change occurs.

\subsection{How a general CM amplitude becomes an operator}
For integer coefficients $n_j$, let
\begin{equation}
 L_{\widetilde A}=n_0I+n_1P+n_2P^2+n_3P^3
\end{equation}
act by convolution on another lifted CM. Put $u=n_0-n_2$ and $v=n_1-n_3$. Iterating \cref{eq:intertwiner} gives
\begin{equation}\label{eq:operatorbridge}
 QL_{\widetilde A}=
 \underbrace{\begin{pmatrix}u&-v\\v&u\end{pmatrix}}_{uI+vJ}\,Q.
\end{equation}
Thus multiplication by a CM-valued phase coefficient is not merely a verbal analogy: it is a four-position convolution operator whose quotient is multiplication by the pair $(u,v)$. On pairs, the induced product is
\begin{equation}
 (u,v)\cdot(s,t)=(us-vt,ut+vs).
\end{equation}
This is the usual Gaussian-integer product in real coordinates. Writing pairs instead of $i$ removes a symbol from the implementation, not the mathematical complex structure.

\subsection{What is lost, and the exact correction for XOR}
The map $Q$ is not injective. Among the sixteen binary arrays its image contains nine pairs in $\{-1,0,1\}^2$. The zero CM, $I_2$, the XOR truth-table matrix, and the all-ones CM all map to $(0,0)$. These are different predicates. Therefore \emph{one must not use $Q$ as a semantics-preserving compression of arbitrary truth-table CMs}.

Nor does $Q$ preserve the original Boolean operations. For example, $E_0\oplus E_0=0$ while $QE_0+QE_0=(2,0)$. Similarly $E_0\land E_1=0$, whereas $(QE_0)(QE_1)=(0,1)$. The exact additive defect is useful:
\begin{proposition}[Carry/overlap correction]\label{prop:carry}
For binary CMs embedded coefficientwise in the integer lift,
\begin{equation}\label{eq:carry}
 Q(A\oplus B)=QA+QB-2Q(A\land B).
\end{equation}
\end{proposition}
\begin{proof}
Apply the integer identity $a\oplus b=a+b-2ab$ at each coordinate, then the integer-linear map $Q$.
\end{proof}
This formula is one practical way to detect an invalid ``Boolean-to-amplitude'' rewrite. Dropping its overlap term changes the semantics.

Transpose satisfies $Q(A^T)=\overline{QA}$, where pair conjugation sends $(u,v)$ to $(u,-v)$. For a binary array, complement happens to satisfy $Q(\neg A)=-QA$ because $Q(\ones)=0$. That last identity is a consequence of the information-losing quotient; it does not identify Boolean NOT with complex conjugation or with a general physical phase operation.

\section{The signed phase norm and visible CM gate operations}\label{sec:normlift}
\subsection{A quotient norm, not a new Boolean Born rule}
The squared length of the quotient coordinates is
\begin{equation}\label{eq:normQ}
 N_Q(n)=\|Qn\|^2=(n_0-n_2)^2+(n_1-n_3)^2
       =n^T(I-P^2)n.
\end{equation}
The matrix $Q^TQ=I-P^2$ is positive semidefinite with rank two. It is a seminorm on the four raw count coordinates and a positive norm on their quotient. In terms of integer self-correlation,
\begin{equation}
 N_Q(n)=C_{n,n}(0)-C_{n,n}(2).
\end{equation}
Unlike Hamming weight, it regards opposite contributions as cancelling. For example, $N_Q(E_0+E_2)=0$ but $N_Q(2E_0)=4$. For $0/1$ entries it equals
\begin{equation}
 \wt(A)-2(a_0a_2+a_1a_3),
\end{equation}
not $\wt(A)$ in general.

\begin{proposition}[Conditional uniqueness of the phase quadratic form]
Let a real positive-semidefinite quadratic form on $\R^4$ be invariant under $P$, vanish on $\ker Q$, and assign squared length one to $E_0$. Then it is exactly $N_Q$.
\end{proposition}
\begin{proof}
Positive semidefiniteness implies that null vectors lie in the kernel of the representing symmetric matrix, so the form descends to $\R^4/\ker Q\cong\R^2$. Its symmetric matrix $G$ must satisfy $J^TGJ=G$. Writing out a symmetric two-by-two $G$ forces equal diagonal entries and zero off-diagonal entries: $G=\lambda I$. The value at $QE_0$ fixes $\lambda=1$.
\end{proof}
The theorem has substantive assumptions. In particular, specifying $\ker Q$ as physically null selects the primitive quarter-phase sector. Rotation invariance alone also permits contributions from the trivial and sign sectors of the regular representation. In Fourier language, $Q$ selects evaluation at $g=i$ (with its conjugate); the full coefficient squared length instead averages the energies of all four characters. Changing from one form to the other is not merely remembering more phase data. It is changing which directions carry measurable amplitude.

For a collection of branch amplitudes $z_x=Qn_x$, define
\begin{equation}\label{eq:stateC4}
 \Psi=2^{-h/2}\sum_x z_x\ket x,\qquad
 \|\Psi\|^2=2^{-h}\sum_x|z_x|^2.
\end{equation}
If we identify this representation with the usual Hilbert-space quantum model, its outcome probabilities are the usual squared moduli. This is a \emph{pullback of the Born rule under an explicit representation}, not a derivation of that empirical probability law from AND, XOR, or rotational symmetry alone.

\subsection{An \texorpdfstring{$S$}{S} gate that really rotates a CM amplitude}
Each logical branch carries an entire lifted CM, rather than one column of a shared CM. In the branch order $(\ket0,\ket1)$, define
\begin{equation}
 S_{\rm CM}(A,B)=(A,PB).
\end{equation}
Applying $Q$ separately to both branches gives
\begin{equation}
 (QA,QB)\longmapsto(QA,JQB).
\end{equation}
A common rotation of every quotient branch multiplies the whole state by a global phase and does not change its physical predictions. Only relative rotations between branches are observable; four phase labels are not four distinguishable physical states of a lone occupied basis branch.

This is the exact operator meaning of multiplying the $\ket1$ amplitude by a quarter phase. The controlled choice concerns the \emph{branch label}; the rotation acts on all four phase coordinates of its coefficient. This resolves the earlier masked-column confusion.

\subsection{A CM-native splitter and recombiner}
One can keep the geometry visible even for Hadamard. Define on raw counts
\begin{equation}\label{eq:CMH}
 H_{\rm CM}:(A,B;h)\longmapsto(A+B,\ A+P^2B;\ h+1).
\end{equation}
The $+$ signs here count paths over $\Z$, not modulo two. By $QP^2=-Q$, its quotient is
\begin{equation}
 H:(u,v;h)\longmapsto(u+v,u-v;h+1).
\end{equation}
It preserves the normalized quotient norm because
\begin{equation}
 |u+v|^2+|u-v|^2=2(|u|^2+|v|^2).
\end{equation}
The raw four-phase map is not required to be invertible on directions annihilated by $Q$; the induced map on the quotient is unitary. An alternative raw representation uses signed subtraction directly. The two are equivalent after $Q$, not as operations on original Boolean truth tables.

\paragraph{Worked example with every CM visible.}
Start with $E_0\ket0$. Splitting gives $(E_0\ket0+E_0\ket1)/\sqrt2$. A quarter-turn on the second branch gives
\begin{equation}
 \frac{E_0\ket0+E_1\ket1}{\sqrt2}.
\end{equation}
Recombination by \cref{eq:CMH} gives raw coefficients $E_0+E_1$ and $E_0+E_3$, with denominator two. Their quotient coordinates are $(1,1)$ and $(1,-1)$, with squared lengths two each. Thus both outcomes have probability $1/2$ in the quantum interpretation.

With no phase shift, the raw outputs are $2E_0$ and $E_0+E_2$. Their quotient lengths are four and zero. With a half-turn the outputs reverse. This is the precise distinction that XOR could not supply: equal directions reinforce by retaining multiplicity, while opposite directions cancel under the quotient.

\subsection{Tensor products and the location of entanglement}
For independent registers, tensor their branch labels and multiply their amplitude coefficients using convolution before quotienting, or Gaussian multiplication afterward. The compatibility follows from $Q(A\star B)=QA\,QB$ in the \emph{integer} convolution algebra. Consequently
\begin{equation}
 \|\Psi\otimes\Phi\|^2=\|\Psi\|^2\|\Phi\|^2.
\end{equation}
CNOT sends $\ket{x,y}$ to $\ket{x,x\oplus y}$ and simply permutes the carried coefficients. Applying a splitter and CNOT to $\ket{00}$ yields $(\ket{00}+\ket{11})/\sqrt2$ in the quotient representation. This is an exact encoding of a standard entangled state, not evidence that a classical array implementation has become a nonlocal physical system.

A previous strict-ring state, $[\neg R]\ket{11}+[R]\ket{00}$, was also nonfactorable over $\alg_2$: all $16^4$ possible two-factor assignments failed. Its coefficient matrix has determinant $[\neg R]\star[R]=1+g^2\ne0$, whereas an outer product over a commutative ring has zero determinant. This algebraic nonseparability is correctly distinguished from the Hilbert-space and experimental senses of entanglement.
